\documentclass[12pt,a4paper]{article}

\usepackage{fontspec}
\usepackage{polyglossia}
\setmainlanguage{russian}
\setotherlanguage{english}
\setmainfont{IBM Plex Serif}
\setmonofont{IBM Plex Mono}[Scale=0.75]
\usepackage{geometry}
\usepackage{hyperref}
\usepackage{listings}
\usepackage{xcolor}
\usepackage{amsmath}
\usepackage{amssymb}
\usepackage{booktabs}
\usepackage{array}
\usepackage{enumitem}
\usepackage{fancyhdr}
\usepackage{titlesec}

\geometry{
  left=1cm, right=1cm,
  top=2cm, bottom=2cm
}

\begin{document}

\title{Декодер DEFLATE (RFC 1951)}
\author{Головин Иван, Б05-304}
\date{}

\maketitle

% ============================================================
\section{Постановка задачи}

\subsection{Неформальная постановка}

Надо написать программу, декодирующую данные, сжатые алгоритмом DEFLATE. Реализация выполняется на языке Haskell
без внешних библиотек.

\subsection{Формальная постановка}

Пусть входные данные задаются строкой в кодировке Base64, которая передаётся
единственным аргументом командной строки. Программа должна выводить декодированные
данные в стандартный вывод.

\subsubsection{Входная строка Base64}

Алфавит Base64: $\Sigma_{64} = \{\,\text{A}\dots\text{Z},\ \text{a}\dots\text{z},\ \text{0}\dots\text{9},\ \text{+},\ \text{/}\,\}$.

\[
\begin{aligned}
Base64String &\ \to\  Group^*\  FinalGroup \\
Group        &\ \to\  c_1\,c_2\,c_3\,c_4,\quad c_i \in \Sigma_{64} \\
FinalGroup   &\ \to\  c_1\,c_2\,c_3\,c_4 \ \mid\  c_1\,c_2\,c_3\,\text{=} \ \mid\  c_1\,c_2\,\text{==}
\end{aligned}
\]

Каждые четыре символа кодируют три байта (24 бита). Символ '=' является
заполнителем и означает, что последняя группа кодирует меньше байт.

\subsubsection{Поток байт DEFLATE}

Согласно RFC\,1951 поток байт DEFLATE представляет собой последовательность
блоков:

\[
DeflateStream \ \to\  Block^+
\]

Каждый блок начинается заголовком из трёх бит:

\begin{center}
\begin{tabular}{lll}
\toprule
Поле & Биты & Значение \\
\midrule
\textsc{bfinal} & 1 & 1 если это последний блок \\
\textsc{btype}  & 2 & тип блока (00, 01, 10) \\
\bottomrule
\end{tabular}
\end{center}

Типы блоков:
\[
Block \ \to\
  \textsc{bfinal}\ \textsc{btype}\
  \begin{cases}
    StoredBlock  & \text{если } \textsc{btype} = 00_2 \\
    FixedBlock   & \text{если } \textsc{btype} = 01_2 \\
    DynamicBlock & \text{если } \textsc{btype} = 10_2
  \end{cases}
\]

\paragraph{Блок без сжатия ($\textsc{btype} = 00_2$):}
\[
StoredBlock \ \to\
  \underbrace{padding}_{\text{выравнивание до байта}}\
  \underbrace{LEN}_{\text{2 байта}}\
  \underbrace{NLEN}_{\text{дополнение}}\
  \underbrace{byte^{LEN}}_{\text{данные}}
\]

\paragraph{Блок с фиксированными кодами ($\textsc{btype} = 01_2$):}
\[
FixedBlock \ \to\  Symbol^*\  EOB
\]
где таблицы кодов Хаффмана фиксированы стандартом (раздел~\ref{sec:fixed}).

\paragraph{Блок с динамическими кодами ($\textsc{btype} = 10_2$):}
\[
DynamicBlock \ \to\
  \underbrace{HLIT}_{\text{5 бит}}\
  \underbrace{HDIST}_{\text{5 бит}}\
  \underbrace{HCLEN}_{\text{4 бита}}\
  \underbrace{CodeLenCodes}_{\text{HCLEN+4 групп по 3 бита}}\
  \underbrace{LitLenCodes}_{\text{HLIT+257 значений}}\
  \underbrace{DistCodes}_{\text{HDIST+1 значений}}\
  Symbol^*\  EOB
\]

\paragraph{Символы (LZ77-пары):}
\[
Symbol \ \to\
  \begin{cases}
    literal                   & \text{код } 0\dots255\text{: байт данных} \\
    EOB                       & \text{код } 256\text{: конец блока} \\
    LenCode\ DistCode& \text{код } 257\dots285\text{: пара (длина, дистанция)}
  \end{cases}
\]

\subsubsection{Задача декодирования}

Пусть $f_{b64}: \Sigma_{64}^* \to \{0,\dots,255\}^*$ — функция декодирования Base64,
$f_{def}: \{0,\dots,255\}^* \to \{0,\dots,255\}^*$ — функция декодирования DEFLATE.

Требуется реализовать функцию:
\[
\boxed{
  F(s) \ =\  f_{def}\!\left( f_{b64}(s) \right),
  \quad s \in \Sigma_{64}^*
}
\]
и вывести $F(s)$, интерпретируя байты как символы UTF-8.

% ============================================================
\section{Используемые алгоритмы}

\subsection{Декодирование Base64}

Base64 кодирует каждые 3 байта (24 бита) четыремя символами из алфавита $\Sigma_{64}$, потому что каждый символ несёт 6 бит информации ($\log_2 64 = 6$).

\paragraph{Алгоритм.}
\begin{enumerate}
  \item Разбить строку на группы по 4 символа.
  \item Для каждого символа $c$ найти его номер $v(c) \in [0,63]$.
  \item Склеить 4 числа по 6 бит в 24-битное целое:
        $n = v(c_1) \cdot 2^{18} + v(c_2) \cdot 2^{12} + v(c_3) \cdot 2^6 + v(c_4)$.
  \item Извлечь три байта: $b_1 = n \gg 16$,\  $b_2 = (n \gg 8)\ \&\ \text{0xFF}$,\  $b_3 = n\ \&\ \text{0xFF}$.
  \item Если последняя группа оканчивается на '=': одним '=' - выдать 2 байта;
        двумя - выдать 1 байт.
\end{enumerate}

\subsection{Монадический парсер битового потока}

Всё декодирование построено на типе:
\[
\text{type BitParser a = StateT BitState Maybe a}
\]

где \text{BitState} хранит список оставшихся байт и позицию следующего бита
внутри текущего байта (0 = LSB). Использование монады \text{StateT} позволяет
передавать состояние потока неявно через монадическую цепочку, а монада
\text{Maybe} обеспечивает "отказ" при несоответствии входных данных.

\paragraph{Порядок битов в DEFLATE.}

Биты внутри байта читаются от младшего к старшему (LSB first).
Коды Хаффмана при этом интерпретируются в обратном порядке:
первый прочитанный бит является старшим битом кода.

\subsection{Коды Хаффмана}
\label{sec:huffman}

\subsubsection{Построение дерева по длинам кодов}

RFC 1951 задаёт таблицы Хаффмана не явным перечислением кодов,
а списком длин кодов для каждого символа. Дерево восстанавливается
по следующему алгоритму:

\begin{enumerate}
  \item Подсчитать $bl\_count[l]$ — количество кодов длины $l$.
  \item Вычислить начальный код для каждой длины:
        \[
          next\_code[1] = 0;\quad
          next\_code[l+1] = \left(next\_code[l] + bl\_count[l]\right) \cdot 2.
        \]
  \item Присвоить каждому символу с длиной $l$ очередной код $next\_code[l]$,
        после чего увеличить $next\_code[l]$ на 1.
  \item Вставить тройки $(\text{ символ}, \text{ длина}, \text{ код})$
        в бинарное дерево, спускаясь по битам кода от старшего к младшему:
        бит 0 — в левого ребёнка, бит 1 — в правого.
\end{enumerate}

\subsubsection{Декодирование символа}

Начиная с корня дерева, читать по одному биту и переходить
в левого (бит=0) или правого (бит=1) ребёнка. Достигнув листа,
вернуть хранящееся там значение.

\subsubsection{Фиксированные таблицы}
\label{sec:fixed}

При $\textsc{btype}=01_2$ таблица задана стандартом:

\begin{center}
\begin{tabular}{lll}
\toprule
Символы & Длина кода & Примечание \\
\midrule
0--143   & 8 бит & литералы (ASCII + расширенный) \\
144--255 & 9 бит & литералы выше 127 \\
256--279 & 7 бит & EOB + коды длин \\
280--287 & 8 бит & коды длин \\
\bottomrule
\end{tabular}
\end{center}

Дистанции кодируются 5-битными кодами (30 значений, длина 5 для всех).

\subsection{Алгоритм LZ77 (декодирование)}

DEFLATE использует LZ77-сжатие: повторяющиеся фрагменты заменяются
ссылкой (длина, дистанция). При декодировании надо:

\begin{enumerate}
  \item Прочитать символ из таблицы литералов/длин.
  \item Если символ $< 256$: добавить байт в выходной буфер.
  \item Если символ $= 256$: конец блока.
  \item Если символ $\in [257, 285]$:
    \begin{enumerate}
      \item Вычислить длину: $L = base\_len[code-257] + \text{extra bits}$.
      \item Прочитать код дистанции из дерева дистанций.
      \item Вычислить дистанцию: $D = base\_dist[dist\_code] + \text{extra bits}$.
      \item Скопировать $L$ байт из позиции $(len(buf) - D)$ выходного буфера.
            Копирование может перекрываться (если $L > D$), что реализует RLE.
    \end{enumerate}
\end{enumerate}

Таблицы кодов длин и дистанций в коде называются, соответственно, $lengthTable$ и $distTable$

\subsection{Декодирование динамических таблиц Хаффмана}

При $\textsc{btype}=10_2$ таблицы передаются в самом потоке,
что позволяет адаптировать коды под конкретные данные.

\paragraph{Алгоритм.}
\begin{enumerate}
  \item Прочитать $HLIT$ (5 бит), $HDIST$ (5 бит), $HCLEN$ (4 бита).
        Количество кодов: литералы = $HLIT+257$, дистанции = $HDIST+1$,
        коды = $HCLEN+4$.
  \item Прочитать $HCLEN+4$ трёхбитных значения в порядке:\\
        $[16,17,18,0,8,7,9,6,10,5,11,4,12,3,13,2,14,1,15]$.
        Это длины кодов для дерева.
  \item Построить дерево Хаффмана.
  \item С помощью дерева прочитать $(HLIT+257) + (HDIST+1)$
        длин кодов:
        \begin{itemize}
          \item Код $0\dots15$: прямое значение длины.
          \item Код $16$: повторить предыдущую длину $3+\text{(2 бита)}$ раз.
          \item Код $17$: $3+\text{(3 бита)}$ нулей.
          \item Код $18$: $11+\text{(7 бит)}$ нулей.
        \end{itemize}
  \item Построить деревья литералов и дистанций.
\end{enumerate}

% ============================================================
\section{Архитектура программы}

Программа организована как единственный модуль Main со следующими
логическими блоками:

\begin{center}
\begin{tabular}{lp{9cm}}
\toprule
Блок & Ответственность \\
\midrule
\text{BitState} / \text{BitParser} & Тип состояния и монада битового парсера \\
\text{bit}, \text{bits}, \text{rawByte} & Примитивы чтения битов и байт \\
\text{decodeBase64} & Декодирование входной строки Base64 \\
\text{HuffmanTree}, \text{buildFromList}, \text{assignCodes} & Дерево Хаффмана и его построение \\
\text{decodeSymbol} & Чтение одного символа по дереву \\
\text{fixedLitLenTree}, \text{fixedDistTree} & Фиксированные таблицы RFC\,1951 \\
\text{lengthTable}, \text{distTable} & Таблицы длин и дистанций LZ77 \\
\text{decodeLZ} & Основной цикл LZ77-декодирования \\
\text{decodeDynTrees} & Чтение динамических таблиц Хаффмана \\
\text{decodeCodeLengths} & RLE-раскодирование длин кодов \\
\text{decodeStored} & Раскодирование блока без сжатия \\
\text{decodeBlocks} & Цикл по блокам DEFLATE \\
\text{decompressBase64} & Точка входа в декодирование \\
\text{main} & Обработка аргументов командной строки \\
\bottomrule
\end{tabular}
\end{center}

\subsection{Конвейер обработки данных}

\[
\text{String (Base64)}
\ \xrightarrow{f_{b64}}\
\text{[Word8]}
\ \xrightarrow{\text{BitState}}\
\text{BitParser}
\ \xrightarrow{\text{decodeBlocks}}\
\text{[Word8]}
\ \xrightarrow{\text{toEnum}}\
\text{String}
\]

\end{document}
